\section{Moderated Network Models}
Here I outline the details of estimating Moderated Network Models (MNMs; e.g.,) \cite{haslbeck2018moderated}. I was under the impression that these already extend to MGMs, but it does seem that it's the case \cite[33]{haslbeck2018moderated}. Lets see.

\subsection{Basics of Moderated Regression}
$C$ is a moderator if the effect of predictor $B$ on response variable $A$ is a \textit{linear} function of $C$
\begin{align} \label{linMod}
    A &= \beta_B B + \beta_C C + \omega_{BC} BC + \epsilon \\
    &= (\beta_B + \omega_{BC} C) B + \beta_C C + \epsilon \\
    &= \beta_B B + (\beta_C + \omega_{BC} B) C + \epsilon
\end{align}
where the effect of $B$ on $A$ is now equal to $(\beta_B + \omega_{BC} C)$. It's important to mean-center predictors for various reasons \cite[7]{haslbeck2018moderated}. More discussion can be found in \cite{aiken1991} and \cite{afshartous2011}.

The regression framework refers to effects from the conditional distribution framework, while networks speak of effects with reference to the joint distribution. Thus, in regression we talk of $\beta_B$, for instance, as a main effect, or conditional main effect in the presence of moderation, and $\omega_{BC}$ as the moderation/interaction effect. From a network perspective, we refer to $\beta_B$ as a pairwise interaction (it is associated with the product of $AB$ in the joint distribution), and $\omega_{BC}$ as a 3-way interaction (as it is associated with the product $ABC$ in the joint distribution) or a moderation effect (as it moderates the pairwise interaction $AB$). 

\subsection{Construction of the MNM}
The goal is to construct a joint distribution over $p$ continuous variables which allows that each pairwise interaction between variables $X_i$ and $X_j$ is a linear function of all other variables; i.e., the interaction between $X_i$ and $X_j$ is linearly moderated by all other variables. The joint distribution is specified by adding moderation effects to the multivariate Gaussian distribution, which otherwise models (unmoderated) linear pairwise interactions between variables

The multivariate Gaussian distribution can be written as
\begin{equation} \label{multiGauss}
    P(X) = \exp\Bigg\{\sum_i^p \alpha_i \frac{X_i}{\sigma_i} + \mathop{\sum_{i,j \in V}}_{i \neq j} \beta_{i,j} \frac{X_i}{\sigma_i} \frac{X_j}{\sigma_j} + \sum_i^p \frac{X_i^2}{\sigma_i^2} - \Phi(\alpha,\beta)\Bigg\}
\end{equation}
where $X \in \mathbb{R}^p,V=\{1,2,\ldots,p\}$ is the index set for the $p$ variables, $\alpha$ is the $p$-vector of intercepts, $\beta$ is a $p \times p$ matrix of $\binom{p}{2}$ partial correlations.\footnote{$\binom{n}{p} = \frac{n!}{k!(n-k)!}$, the $k^{\text{th}}$ binomial coefficient of the polynomial expansion of $(1+x)^n$} $\sigma_i^2$ is the variance of $X_i$, and $\Phi(\alpha,\beta)$ is the log normalizing constant that ensures the probability distribution integrates to 1.

To extend the distribution in \eqref{multiGauss} with all possible moderation effects, 3-way interactions are added to the model
\begin{equation} \label{multiGauss3}
    \begin{split}
        P(X) = \exp\Bigg\{&\sum_i^p \alpha_i \frac{X_i}{\sigma_i} + \mathop{\sum_{i,j \in V}}_{i \neq j} \beta_{i,j} \frac{X_i}{\sigma_i} \frac{X_j}{\sigma_j} + \\
        &\mathop{\sum_{i,j,q \in V}}_{i \neq j \neq q} \omega_{i,j,q} \frac{X_i}{\sigma_i} \frac{X_j}{\sigma_j} \frac{X_q}{\sigma_q} + \sum_i^p \frac{X_i^2}{\sigma_i^2} - \Phi(\alpha,\beta,\omega)\Bigg\}
    \end{split}
\end{equation}
where $\omega$ is a $p \times p \times p$ array of $\binom{p}{3}$ 3-way interactions (moderation effects). One may also decide to only include a subset of moderation effects, rather than all moderation effects. The distribution of $X_s$ conditioned on all remaining variables $X_{\setminus s}$ is given by
\begin{equation} \label{condGauss3}
    \begin{split}
        P(X_s|X_{\setminus s}) = \exp\Bigg\{&\alpha_s \frac{X_s}{\sigma_s} + \mathop{\sum_{i \in V}}_{i \neq s} \beta_{i,s} \frac{x_i}{\sigma_i} \frac{X_s}{\sigma_s} + \\
        &\mathop{\sum_{i,j \in V}}_{i \neq j \neq s} \omega_{i,j,s} \frac{x_i}{\sigma_i} \frac{x_j}{\sigma_j} \frac{X_s}{\sigma_s} + \frac{X_s^2}{\sigma_s^2} - \Phi^*(\alpha,\beta,\omega) \Bigg\}
    \end{split}
\end{equation}
where $\Phi^*(\alpha,\beta\omega)$ is the log-normalizing constant with respect to the conditional distribution $P(X_s|X_{\setminus s})$, and lowercase letters $x_i$ indicate fixed values as opposed to random variables $X_i$. Below it is shown that \eqref{condGauss3} is a conditional Gaussian distribution, where $\sigma_s = 1$ is used for simplicity
\begin{gather} \label{mu3}
    \mu_s = \alpha_s + \summ{i \in V}{i \neq s} \beta_{i,s} x_i + \summ{i,j \in V}{i \neq j \neq s} \omega_{i,j,s} x_i x_j \\
    \Phi^*(\alpha,\beta,\omega) = \frac{\mu_s}{2} - \log \Big(\frac{1}{\sqrt{2\pi}}\Big)
\end{gather}
from this we can rewrite \eqref{condGauss3} into
\begin{equation}
    \begin{split}
        P(X_s|X_{\setminus s}) &= \frac{1}{\sqrt{2\pi}} \exp\Bigg\{\mu_s X_s - \frac{X_s}{2} - \frac{\mu_s}{2}\Bigg\} \\
        &= \frac{1}{\sqrt{2\pi}} \exp\Bigg\{\frac{(X_s - \mu_s)^2}{2}\Bigg\}
    \end{split}
\end{equation}
which is the standard form of the conditional Gaussian distribution.

Since each $P(X_s|X_{\setminus s})$ is only parameterized by $\mu_s$ and regression estimates $\mathbb{E} \lbrack X_s|X_{\setminus s} \rbrack = \mu_s$, the above derivation shows that a MNM can be estimated through a series of $p$ regressions that include moderation effects. I.e., the equation for the mean $\mu_s$ of the conditional distribution $X_s$ in \eqref{mu3} has the same form as the moderated linear regression in \eqref{linMod} except with more terms.

The mean \eqref{mu3} of the conditional distribution $P(X_s|X_{\setminus s})$ compared to the joint distribution $P(X)$ in \eqref{multiGauss3} explains the different terminology for interaction effects. For instance, in the joint distribution the second term indicates pairwise interactions because two variables are multiplied. In the mean of the conditional, the second term only contains a single variable and is thus referred to as a main effect. 

Given the estimation of $\ell_1$-regularized regressions, however, the fact that 3-way interactions are included in the model does not necessitate that all relevant pairwise interactions must also be included \cite[9--10]{haslbeck2018moderated}. It is possible for the LASSO procedure to set the coefficients for some pairwise interactions to zero, even if 3-way interactions moderating that pairwise interaction are nonzero. This may be considered problematic, and may be a reason to not use the $\ell_1$-regularization estimation procedure. However, we may adapt this constraint to take into account such relationships \cite{bien2013lasso}.

While the MNM joint distribution \eqref{multiGauss3} can be factorized into $p$ conditional Gaussian distributions, it is not a multivariate Gaussian distribution. For the multivariate Gaussian distribution, the covariance matrix must be positive-definite to ensure that the distribution is normalizable. However, the constraints for normalizability for this model \eqref{multiGauss3} are unknown \cite{yang2014mixed}. The consequence is that, for an estimated MNM, we do not know whether the probabilities computed under the model are actual probabilities. Nevertheless, one can make probabilistic statements and predictions from the conditional distributions without limitations. Given that most network analyses in psychological research do not perform inference on the joint distribution but analyze the relative size of parameters and compute predictions from conditional distributions, the practical implications of this limitation are small \cite[10]{haslbeck2018moderated}.

To visualize the MNM, each interaction can be represented via \textit{factor nodes} \cite{koller2009}. We can distinguish between \textit{partial} and \textit{full} moderation, respectively, as: (a) when variables $A$ and $B$ have a pairwise interaction $AB$, but also interact with a third variable $C$, $ABC$, and (b) when variables $A$ and $B$ do \textit{not} have a pairwise interaction (i.e., $AB$ is estimated at 0 via the $\ell_1$-regularized regressions), yet interact with a third variable $C$, $ABC$.

\subsection{Estimation via Nodewise Regression}
This approach is similar to estimating the multivariate Gaussian distribution \cite{meinshausen2006high}, except here the conditionals of the MNM \eqref{condGauss3} are estimated rather than the joint Gaussian. 

To estimate $p$ regression models, we minimize the squared loss plus the $\ell_1$-norm of the parameter vectors $\beta_{s,\cdot}$ and $\omega_{s,\cdot,\cdot}$ for each variable $s \in V$
\begin{equation} \label{lassoMNM}
    \argmin_{\beta_{s,\cdot},\omega_{s,\cdot,\cdot}} \Bigg\{\sum_{z=1}^n (X_{z,s} - \hat{\mu}_s)^2 + \lambda_s \left(\norm{\beta_{s,\cdot}}_1 + \norm{\omega_{s,\cdot,\cdot}}_1\right)\Bigg\}
\end{equation}
where $X_{z,s}$ is the value of variable $s$ in row $z$ of the data matrix, $\hat{\mu}_{z,s}$ is predicted mean for row $z$ \eqref{mu3}, $\beta_{s,\cdot}$ and $\omega_{s,\cdot,\cdot}$ are vectors containing parameters associated with pairwise and 3-way interactions, respectively. $\norm{\beta_s}_1 + \norm{\omega_s}_1 = \summ{i \in V}{i \neq s} |\beta_{s,i}| + \summ{i,j \in V}{i \neq j \neq s} |\beta_{s,i,j}|$ is the sum of the $\ell_1$-norms of both parameter vectors, and $\lambda_s$ is the tuning parameter that weights the $\ell_1$-norm relative to the squared loss. 

Prior to estimation, all variables are mean-centered and divided by their standard deviation. This ensures that the penalization parameter doesn't depend on the standard deviation of its associated (product of) variables and also simplifies the model because all intercepts are equal to zero. Mean-centering is necessary to obtain interpretable parameter estimates.

There are three reasons for choosing to use $\ell_1$-regularized (LASSO) regression: (a) the number of parameters can be large if $p$ is large and many moderators are included in the model, which leads to high variance in the parameter estimates (overfitting); (b) given that $\ell_1$-regularization shrinks small parameter estimates to zero, this helps deal with potentially large $p$ data, as well as with interpretation from a graph-theoretic perspective; and (c) this ensures that the model is identifiable when there are more parameters than observations. For example, if we have an MNM with $p=20$ variables, including moderating effects, then each nodewise regression has $p-1 + \binom{p-1}{2}=190$ parameters, meaning that unregularized methods would require $n \geq 190$ observations.

The main assumption underlying $\ell_1$-regularized regression is that most of the parameters in the true model are equal to zero (this is called the \textit{sparsity} assumption). This seems reasonable for MNMs with psychological data, as we would not expect that every variable moderates every pairwise interaction. See \cite{hastie2015statistical} for a detailed discussion of $\ell_1$-regularized regression. Moreover, $\lambda_s$ must be selected to tune the strength of penalization. When $\lambda_s = 0$, the loss function \eqref{lassoMNM} reduces to the standard squared loss as in OLS regression. Cross-validation can be used to select the optimal $\lambda_s$, or a model-selection criterion such as the EBIC, which has been shown to perform well in estimating sparse parameter vectors \cite{foygel2010extended} better than the BIC \cite{schwarz1978estimating}.

\subsubsection{Combine Estimates to Joint MNM}
The above estimation produces two estimates for every pairwise interaction effect (e.g., $\beta_{(i,j)}$ for the regression of $X_j$ on $X_i$, and $\beta_{(j,i)}$ for the regression of $X_i$ on $X_j$), and three estimates for every 3-way interaction: (a) $X_q$ moderating the predictor $X_j$ in the regression on $X_s$, (b) $X_j$ moderating the predictor $X_q$ in the regression on $X_s$, and (c) $X_j$ moderating the predictor $X_s$ in the regression on $X_q$. To aggregate the coefficients, we either take the arithmetic mean across the two/three values (OR-rule), or take the arithmetic mean across the two/three values if all three values are nonzero and otherwise set the aggregated parameter to zero (AND-rule). 

\subsection{Limitations}
\begin{enumerate}
    \item There is no explicit constraint on the parameter space of the MNM joint distribution that ensures that the distribution is normalizable. This can be worked around by using a rejection sampler, if the goal is to simulate data from an MNM joint distribution. However, for a given model estimated from the data, it is unclear whether the joint distribution is normalizable. Thus, while the joint distribution does correctly capture the dependency structure in the data, it might not be possible to define a probability distribution over it. Consequently, when attempting perform inference on the model we are not guaranteed to obtain probabilities. However, all conditional distributions \textit{are} proper distributions; thus, there shouldn't be an issue with making inferences from the conditional distributions, and predictions for any variable can be computed without limitation.
    \item Performance in recovering true parameters may differ depending on the data and how the model is specified. 
    \item May be challenging to recover moderation effects with $\ell_1$-regularization, given that these tend to be small. But this will be an issue with any estimation algorithm. Could use a more liberal tuning parameter $(\gamma = 0)$, but this would increase the false positive rates for all variables.
\end{enumerate}
We can think of interaction effects simply as additional variables in a multiple regression, and thus the performance in recovering parameters is in principle the same as it is for main effects. Therefore one can draw on the rich literature on these models to make predictions about the performance in recovering MNMs for different setups \cite{hastie2015statistical,buhlmann2011statistics}.

\subsection{Future Directions}
\begin{enumerate}
    \item $\ell_1$-regularized nodewise regression was used to estimate MNMs, and was chosen because it deals well with large $p$ and renders interpretation easier by setting small parameters to zero. However, some situations (e.g., correlated predictors in skewed data) may be better dealt with by other estimators (e.g., significance testing).
    \item Since all parameters are freely estimated, it is possible that a moderation effect is estimated as nonzero even when all pairwise interactions that it's based on are estimated at zero. One may enforce a \textit{structural hierarchy} to avoid this issue \cite{bien2013lasso}, but this will reduce the fit of the model.
    \item Instead of using a single regularization parameter $\lambda_s$ for the entire parameter vector (which includes both pairwise and 3-way interactions), we may assume a different level of sparsity for pairwise and moderation effects and use separate regularization parameters. The downsides are that this increases computational complexity (the model selection procedure searches a grid of $\lambda_s$ rather than a sequence), and estimates may become more unstable since the considered model space is larger.
    \item Would be useful to extend the present work to include MGMs. In principle the same approach here could be used there: add linear moderation effects for one or several moderator variables. The only difference is that for MGMs we perform nodewise regressions using the GLM framework with appropriate link-functions \cite{nelder1972generalized}. Would be useful to detail the possible moderated interaction types between different variables, describe how to interpret them, and provide performance results for estimating different types of moderation effects in different situations. For example, an MGM with a single categorical variable could allow one to estimate group differences across two or several groups in a principled way that is likely to out-perform group-split based methods.
    \item It would be useful to extend the notion of centrality to factor graphs and obtain \textit{moderator centrality}. A na{\"i}ve way to do this would be to simply add up the moderation effect of a given variable to find out which one has the strongest influence on the pairwise interactions in the model. Could be especially useful in data with contextual variables, as it could allow one to identify which of them have the strongest influence on a network model of psychological variables or symptoms.
\end{enumerate}
In sum, the Moderated Network Model relaxes the assumption of Gaussian Graphical Models that each pairwise interaction is independent of the value of all other variables by allowing that each pairwise interaction is \textit{moderated} by (potentially) all other variables. This allows more precise statements about the sign and value of a given interaction parameter in a given situation, and may allow us to estimate more accurate models for subgroups and individuals.